\chapter{Gastritis}
\index{Gastritis}

\section*{Definition and Overview}
Gastritis is inflammation of the lining of the stomach. It can occur suddenly (acute gastritis) or develop gradually and persist for a long time (chronic gastritis), and it ranges from a mild irritation that causes no symptoms to erosions and bleeding. The most common causes are infection with the bacterium Helicobacter pylori, long-term use of non-steroidal anti-inflammatory drugs (NSAIDs) such as ibuprofen, and excessive alcohol. Symptoms may include upper abdominal pain, nausea, bloating, and loss of appetite, though many people have no symptoms at all. Gastritis usually resolves with treatment of the cause and medicines that reduce stomach acid. This chapter explains the causes, symptoms, and treatment of gastritis.

\section*{Epidemiology}
Gastritis is common, and H.\ pylori infection is its most frequent cause, affecting around a quarter to a third of people, with higher rates in older people and in some migrant populations. Around half of the world's population carries H.\ pylori, though most have no symptoms. NSAID use is widespread, and long-term users are at higher risk of gastritis and ulcers. Chronic gastritis from H.\ pylori can progress to peptic ulcers and, over decades, to stomach cancer in a small minority. Acute gastritis from alcohol and NSAIDs is very common and usually resolves when the cause is removed.

\section*{Aetiology and Pathophysiology}
\begin{itemize}
    \item \textbf{H.\ pylori:} the bacterium colonises the stomach lining, causing chronic inflammation that can persist for life without treatment
    \item \textbf{NSAIDs:} medicines such as ibuprofen, naproxen, and aspirin reduce the prostaglandins that protect the stomach lining, making it vulnerable to acid
    \item \textbf{Alcohol:} irritates and erodes the stomach lining, especially in large amounts
    \item \textbf{Stress and serious illness:} severe illness, burns, and major surgery can cause stress-related gastritis
    \item \textbf{Other causes:} bile reflux, viral infections, and autoimmune gastritis, in which the immune system attacks the stomach lining
    \item \textbf{The mechanism:} the balance between acid and the protective mucus layer is disturbed, so acid damages the lining, causing inflammation and sometimes erosions
    \item \textbf{Chronic gastritis:} long-standing inflammation can thin the lining, reduce acid production, and increase cancer risk in a minority
    \item \textbf{Complications:} erosions can bleed, and chronic H.\ pylori gastritis can lead to peptic ulcers and stomach cancer
\end{itemize}

\section*{Clinical Features}
\begin{itemize}
    \item Upper abdominal pain or discomfort, often described as burning, gnawing, or aching
    \item Nausea and sometimes vomiting
    \item Bloating, a feeling of fullness, and loss of appetite
    \item Indigestion and belching
    \item No symptoms at all in many people, especially with chronic gastritis
    \item Black, tarry stools or vomiting blood with erosive, bleeding gastritis
    \item Anaemia and fatigue from chronic slow bleeding
    \item Symptoms that worsen with NSAIDs, alcohol, or stress
\end{itemize}

\section*{Differential Diagnosis}
Upper abdominal symptoms have many causes.
\begin{itemize}
    \item Peptic ulcer disease, which is closely related and caused by the same factors
    \item Gastro-oesophageal reflux disease (GORD)
    \item Functional dyspepsia, with similar symptoms but no visible inflammation
    \item Gallbladder disease
    \item Pancreatitis
    \item Stomach cancer, which must be considered with weight loss, bleeding, and older age
    \item Heart attack, which can present as upper abdominal pain
    \item Endoscopy distinguishes these
\end{itemize}

\section*{When to Seek Medical Care}
Anyone with persistent upper abdominal pain, indigestion, nausea, or black stools should be assessed. The combination of age, risk factors, and symptoms guides testing.

\textbf{Contact your local emergency services for:}
\begin{itemize}
    \item Vomiting blood or coffee-ground material
    \item Black, tarry stools, which indicate bleeding from the stomach
    \item Severe abdominal pain
    \item Faintness, dizziness, or collapse, which may indicate significant bleeding
    \item Chest pain or breathlessness
\end{itemize}

\section*{Diagnosis and Investigations}
\begin{itemize}
    \item \textbf{History and examination:} symptoms, medicines, alcohol, and risk factors
    \item \textbf{H.\ pylori testing:} a urea breath test, stool antigen test, or blood antibody test
    \item \textbf{Gastroscopy:} the definitive test, allowing direct inspection of the stomach lining and biopsy
    \item \textbf{Biopsy:} confirms the inflammation and the presence of H.\ pylori or other changes
    \item \textbf{Blood tests:} full blood count for anaemia, and tests of kidney and liver function
    \item \textbf{Review of medicines:} identification of NSAID and other culprit drugs
\end{itemize}

\section*{Management}
\begin{itemize}
    \item \textbf{Acid suppression:} proton pump inhibitors (such as omeprazole) or H2 blockers reduce acid and allow the lining to heal
    \item \textbf{Treat H.\ pylori:} a course of antibiotics with acid suppression ("triple therapy") eradicates the infection
    \item \textbf{Stop or reduce NSAIDs:} use paracetamol where possible, or protect the stomach with a PPI if NSAIDs must continue
    \item \textbf{Reduce alcohol:} limit or stop drinking
    \item \textbf{Stop smoking:} smoking worsens gastritis and delays healing
    \item \textbf{Diet and lifestyle:} eat regular, smaller meals and avoid foods that irritate
    \item \textbf{Treat bleeding:} hospital care with intravenous acid suppression and, if needed, endoscopic treatment
    \item \textbf{Autoimmune gastritis:} monitoring for vitamin B12 deficiency and, long term, specialist surveillance
    \item \textbf{Follow-up:} confirm H.\ pylori eradication and re-test after treatment
\end{itemize}

\section*{Complications}
\begin{itemize}
    \item Peptic ulcers, which can bleed or perforate
    \item Bleeding from erosions, with anaemia or significant haemorrhage
    \item Iron-deficiency anaemia from chronic slow blood loss
    \item Vitamin B12 deficiency in autoimmune gastritis
    \item Stomach cancer, in a small minority with long-standing H.\ pylori gastritis
    \item Recurrence when the cause is not removed
\end{itemize}

\section*{Prognosis and Outcomes}
Gastritis resolves well with treatment. H.\ pylori is eradicated in around 85% of people with standard treatment, healing the inflammation and reducing ulcer and cancer risk. NSAID-related gastritis settles when the medicine is stopped or the stomach is protected. Bleeding gastritis is treated effectively in hospital. The outlook is excellent for most people, and the main tasks are removing the cause and preventing recurrence.

\section*{Prevention and Self-Care}
\begin{itemize}
    \item Use NSAIDs only at the lowest dose for the shortest time, and take them with food
    \item If you need long-term NSAIDs, ask a doctor about stomach protection
    \item Limit alcohol and avoid smoking
    \item Eat regular meals and avoid large, irritating meals late at night
    \item Manage stress
    \item See a doctor for persistent or severe upper abdominal symptoms
    \item Complete the full course of H.\ pylori treatment and confirm eradication
    \item Report black stools or vomiting blood immediately
\end{itemize}

\section*{Sources and Further Reading}
The following reputable sources provide further information for healthcare students and patients seeking reliable, current advice about gastritis.
\begin{itemize}
    \item Healthdirect Australia. \textit{Gastritis.} https://www.healthdirect.gov.au/gastritis
    \item The Gut Foundation. \textit{Gastritis.} https://www.gutfoundation.org.au/
    \item Better Health Channel. \textit{Gastritis.} https://www.betterhealth.vic.gov.au/
    \item NHS. \textit{Gastritis.} https://www.nhs.uk/conditions/gastritis/
    \item Mayo Clinic. \textit{Gastritis.} https://www.mayoclinic.org/
    \item MedlinePlus. \textit{Gastritis.} https://medlineplus.gov/
\end{itemize}
